\documentclass[11pt,a4paper]{article}

\usepackage[T1]{fontenc}
\usepackage[utf8]{inputenc}
\usepackage[ngerman]{babel}
\usepackage{lmodern}
\usepackage{microtype}
\usepackage[a4paper,margin=2.5cm]{geometry}

\usepackage{amsmath,amssymb,amsthm,mathtools}
\usepackage{booktabs}
\usepackage{enumitem}
\usepackage{hyperref}

\hypersetup{
	colorlinks=true,
	linkcolor=blue,
	citecolor=blue,
	urlcolor=blue,
	pdftitle={Ein diskreter Vierlings-Hamiltonian im Bamberger Modell},
	pdfauthor={Thomas Hoffbauer}
}

\title{Ein diskreter Vierlings-Hamiltonian im Bamberger Modell}
\author{Thomas Hoffbauer}
\date{\today}

\newtheorem{definition}{Definition}[section]
\newtheorem{lemma}[definition]{Lemma}
\newtheorem{proposition}[definition]{Proposition}
\newtheorem{remark}[definition]{Bemerkung}
\newtheorem{corollary}[definition]{Korollar}

\newcommand{\SigmaBM}{\Sigma}
\newcommand{\XE}{\mathcal X_E}
\newcommand{\HE}{\mathcal H_E}
\newcommand{\wt}{\mathrm{wt}}

\begin{document}
	\maketitle
	
	\begin{abstract}
		Wir definieren ein diskretes Vierlingsmodell auf dem neutralen Sektor der Klein-Vierergruppe
		\(
		\SigmaBM=\{E,A,B,C\}.
		\)
		Der Zustandsraum besteht aus allen Vierlingen
		\(
		(\sigma_1,\sigma_2,\sigma_3,\sigma_4)\in\SigmaBM^4
		\)
		mit
		\(
		\sigma_1\sigma_2\sigma_3\sigma_4=E.
		\)
		Auf dem zugehörigen komplexen Zustandsraum konstruieren wir einen Hamiltonoperator aus einem zyklischen Shift, klassenerhaltenden Paarflip-Operatoren und einem diagonalen Potential. Das Modell besitzt eine natürliche Viererorbit-Struktur und eine diskrete Holonomie in den Shift-Sektoren. Numerisch zeigt der Grundzustand einen Regimewechsel zwischen zwei konkurrierenden neutralen Ordnungen, nämlich einer ABCE-dominierten voll gemischten Ordnung und einer AABB-dominierten paarig kompensierten Ordnung. Für die Parameterwahl
		\(
		t=1,\ \lambda=0.15,\ \mu=0.4
		\)
		liegt der Umschlag numerisch bei
		\(
		\alpha_c\approx -0.72.
		\)
		Der vollständig triviale Zustand \(EEEE\) bleibt im untersuchten Bereich untergeordnet.
	\end{abstract}
	
	\section{Einleitung}
	
	Ziel dieser Notiz ist die Formulierung eines minimalen diskreten Operatoransatzes für Vierlingskonfigurationen im Rahmen des Bamberger Modells. Im Vordergrund stehen drei Anforderungen:
	
	\begin{enumerate}[label=(\roman*)]
		\item ein endlicher globaler Zustandsraum mit nichttrivialer Neutralitätsbedingung,
		\item eine natürliche Dynamik auf diesem Raum,
		\item eine numerisch überprüfbare Konkurrenz verschiedener globaler Ordnungsformen.
	\end{enumerate}
	
	Der Ansatz ist bewusst elementar. Es wird weder eine Kontinuumsgrenze noch eine physikalische Quantisierung behauptet. Stattdessen wird ein wohldefinierter diskreter Hamiltonoperator untersucht, dessen Spektrum bereits eine nichttriviale Grundzustandsselektion zeigt.
	
	\section{Algebraische Grundlage}
	
	\begin{definition}
		Sei
		\[
		\SigmaBM=\{E,A,B,C\}
		\]
		mit der Struktur der Klein-Vierergruppe versehen:
		\[
		A^2=B^2=C^2=E,\qquad AB=C,\qquad AC=B,\qquad BC=A.
		\]
	\end{definition}
	
	\begin{remark}
		\(\SigmaBM\) ist abelsch, und jedes Element ist sein eigenes Inverses.
	\end{remark}
	
	\begin{definition}
		Der volle Vierlingsraum ist
		\[
		\mathcal X=\SigmaBM^4.
		\]
		Für
		\(
		X=(\sigma_1,\sigma_2,\sigma_3,\sigma_4)\in\mathcal X
		\)
		definieren wir die Gesamtsignatur
		\[
		Q(X):=\sigma_1\sigma_2\sigma_3\sigma_4\in\SigmaBM.
		\]
	\end{definition}
	
	\begin{definition}
		Der neutrale Vierlingssektor ist
		\[
		\XE
		=
		\left\{
		(\sigma_1,\sigma_2,\sigma_3,\sigma_4)\in\SigmaBM^4:
		Q(X)=E
		\right\}.
		\]
		Der zugehörige komplexe Zustandsraum ist
		\[
		\HE=\mathbb C[\XE].
		\]
	\end{definition}
	
	\begin{lemma}
		Es gilt
		\[
		|\XE|=64.
		\]
	\end{lemma}
	
	\begin{proof}
		Drei Einträge können frei gewählt werden; der vierte ist durch die Neutralitätsbedingung eindeutig bestimmt.
	\end{proof}
	
	\section{Operatoren}
	
	\subsection{Shift}
	
	\begin{definition}
		Der Shiftoperator \(T\) auf Basiszuständen ist definiert durch
		\[
		T|\sigma_1,\sigma_2,\sigma_3,\sigma_4\rangle
		=
		|\sigma_2,\sigma_3,\sigma_4,\sigma_1\rangle.
		\]
	\end{definition}
	
	\begin{lemma}
		\(T\) erhält \(\XE\) und erfüllt \(T^4=I\).
	\end{lemma}
	
	\begin{proof}
		Die Erhaltung von \(\XE\) folgt aus der Kommutativität von \(\SigmaBM\). Die Relation \(T^4=I\) ist unmittelbar.
	\end{proof}
	
	\subsection{Paarflips}
	
	\begin{definition}
		Für \(g\in\{A,B,C\}\) und \(1\le i<j\le 4\) definieren wir
		\[
		F_{ij}^{(g)}
		|\sigma_1,\dots,\sigma_i,\dots,\sigma_j,\dots,\sigma_4\rangle
		=
		|\sigma_1,\dots,g\sigma_i,\dots,g\sigma_j,\dots,\sigma_4\rangle.
		\]
	\end{definition}
	
	\begin{lemma}
		Jeder Operator \(F_{ij}^{(g)}\) erhält \(\XE\).
	\end{lemma}
	
	\begin{proof}
		Wegen \(g^2=E\) bleibt das Gesamtprodukt unverändert.
	\end{proof}
	
	\subsection{Potential}
	
	\begin{definition}
		Sei \(V\) ein diagonaler Operator auf \(\HE\), gegeben durch
		\[
		V|X\rangle=v(X)|X\rangle.
		\]
		In den numerischen Rechnungen wird \(v\) aus einem Term \(\mu N_{\mathrm{nt}}(X)\) und einer zusätzlichen Gewichtung der ABCE-Konfigurationen mit Parameter \(\alpha\) zusammengesetzt.
	\end{definition}
	
	\section{Hamiltonoperator}
	
	\begin{definition}
		Für Parameter \(t,\lambda\ge 0\) sei
		\[
		H
		=
		-t(T+T^{-1})
		-\lambda\sum_{1\le i<j\le 4}\sum_{g\in\{A,B,C\}}F_{ij}^{(g)}
		+V.
		\]
	\end{definition}
	
	\begin{lemma}
		\(H\) ist wohldefiniert auf \(\HE\).
	\end{lemma}
	
	\begin{proof}
		Die Invarianz von \(\XE\) unter \(T\) und \(F_{ij}^{(g)}\) wurde bereits gezeigt; \(V\) ist diagonal.
	\end{proof}
	
	\section{Viererorbit und diskrete Holonomie}
	
	Betrachte den Zustand
	\[
	X_0=(A,B,C,E)\in\XE.
	\]
	Sein Orbit unter \(T\) ist
	\[
	X_0,\quad T(X_0),\quad T^2(X_0),\quad T^3(X_0),
	\]
	also explizit
	\[
	(A,B,C,E),\ (B,C,E,A),\ (C,E,A,B),\ (E,A,B,C).
	\]
	
	Auf dem linearen Spann dieses Orbits besitzt \(T\) die Matrixdarstellung
	\[
	T=
	\begin{pmatrix}
		0&0&0&1\\
		1&0&0&0\\
		0&1&0&0\\
		0&0&1&0
	\end{pmatrix}.
	\]
	Die Eigenwerte von \(T\) sind
	\[
	1,\ i,\ -1,\ -i.
	\]
	
	\begin{remark}
		Diese vier diskreten Shift-Sektoren sind die einfachste Form einer Holonomiestruktur im Modell. Der geschlossene Umlauf eines Vierlings ist nicht trivial, sondern trägt eine diskrete Phasenklassifikation.
	\end{remark}
	
	Für den reinen Transportterm
	\[
	H_{\mathrm{trans}}=-t(T+T^{-1})
	\]
	ergeben sich daher auf dem Viererorbit die Eigenwerte
	\[
	-2t,\ 0,\ 0,\ 2t.
	\]
	
	\section{Typklassen neutraler Vierlinge}
	
	\begin{definition}
		Wir unterscheiden drei Typklassen in \(\XE\):
		\begin{enumerate}[label=(\alph*)]
			\item den \emph{EEEE-Typ}
			\[
			(E,E,E,E),
			\]
			\item den \emph{AABB-Typ}, d.\,h. neutrale Vierlinge mit Häufigkeitsmuster \((2,2)\),
			\item den \emph{ABCE-Typ}, d.\,h. neutrale Vierlinge, in denen \(E,A,B,C\) jeweils genau einmal vorkommen.
		\end{enumerate}
	\end{definition}
	
	\begin{remark}
		AABB und ABCE sind zwei verschiedene Formen strukturierter Neutralität. Die numerische Analyse zeigt, dass beide im Grundzustand konkurrieren.
	\end{remark}
	
	\section{Numerische Resultate}
	
	Im Folgenden werden die Parameter
	\[
	t=1,\qquad \lambda=0.15,\qquad \mu=0.4
	\]
	fixiert.
	
	\subsection{Grundzustand bei \texorpdfstring{$\alpha=-0.3$}{alpha=-0.3}}
	
	Für \(\alpha=-0.3\) erhält man numerisch
	\[
	E_0\approx -3.68503172,\qquad E_1\approx -2.22057451.
	\]
	
	Die Typgewichte des Grundzustands sind
	\[
	\wt_{AABB}(E_0)\approx 0.37584705,\qquad
	\wt_{ABCE}(E_0)\approx 0.17788503,\qquad
	\wt_{EEEE}(E_0)\approx 0.05361052.
	\]
	
	Dagegen gilt für den ersten angeregten Zustand
	\[
	\wt_{EEEE}(E_1)\approx 0.67243007.
	\]
	
	\begin{proposition}
		Für \(\alpha=-0.3\) ist der Grundzustand nicht trivial. Er wird von strukturierten neutralen Konfigurationen dominiert, insbesondere vom AABB-Typ, mit signifikanter Beimischung des ABCE-Typs.
	\end{proposition}
	
	\begin{proof}
		Dies folgt unmittelbar aus den angegebenen Typgewichten.
	\end{proof}
	
	\subsection{Scan in \texorpdfstring{$\alpha$}{alpha}}
	
	\begin{table}[h]
		\centering
		\begin{tabular}{rcccc}
			\toprule
			$\alpha$ & $E_0$ & AABB & ABCE & EEEE \\
			\midrule
			$-1.00$ & $-4.00471086$ & $0.427823$ & $0.526006$ & $0.029035$ \\
			$-0.75$ & $-3.87939297$ & $0.467105$ & $0.476415$ & $0.036655$ \\
			$-0.50$ & $-3.76650374$ & $0.505010$ & $0.426820$ & $0.045587$ \\
			$-0.25$ & $-3.66586282$ & $0.540303$ & $0.378661$ & $0.055726$ \\
			$0.00$  & $-3.57694646$ & $0.572016$ & $0.333212$ & $0.066867$ \\
			$0.25$  & $-3.49895031$ & $0.599557$ & $0.291426$ & $0.078731$ \\
			$0.50$  & $-3.43088069$ & $0.622728$ & $0.253861$ & $0.091008$ \\
			\bottomrule
		\end{tabular}
		\caption{Typgewichte des Grundzustands für verschiedene Werte von \(\alpha\).}
	\end{table}
	
	Ein Feinscan im Umschaltbereich ergibt:
	
	\begin{table}[h]
		\centering
		\begin{tabular}{rccccc}
			\toprule
			$\alpha$ & $E_0$ & AABB & ABCE & EEEE & ABCE$-$AABB \\
			\midrule
			$-0.80$ & $-3.90346289$ & $0.459317$ & $0.486381$ & $0.035024$ & $+0.027064$ \\
			$-0.70$ & $-3.85582136$ & $0.474838$ & $0.466450$ & $0.038338$ & $-0.008389$ \\
			$-0.60$ & $-3.81017142$ & $0.490097$ & $0.446565$ & $0.041861$ & $-0.043532$ \\
			$-0.50$ & $-3.76650374$ & $0.505010$ & $0.426820$ & $0.045587$ & $-0.078190$ \\
			$-0.40$ & $-3.72479970$ & $0.519502$ & $0.407307$ & $0.049507$ & $-0.112195$ \\
			$-0.30$ & $-3.68503172$ & $0.533505$ & $0.388113$ & $0.053611$ & $-0.145392$ \\
			$-0.20$ & $-3.64716380$ & $0.546957$ & $0.369319$ & $0.057883$ & $-0.177638$ \\
			$-0.10$ & $-3.61115223$ & $0.559808$ & $0.350997$ & $0.062307$ & $-0.208811$ \\
			$0.00$  & $-3.57694646$ & $0.572016$ & $0.333212$ & $0.066867$ & $-0.238804$ \\
			$0.10$  & $-3.54449005$ & $0.583553$ & $0.316018$ & $0.071543$ & $-0.267535$ \\
			\bottomrule
		\end{tabular}
		\caption{Feinscan des Umschaltbereichs.}
	\end{table}
	
	\begin{proposition}[Numerischer Ordnungswechsel]
		Für die Parameter
		\[
		t=1,\qquad \lambda=0.15,\qquad \mu=0.4
		\]
		zeigt der Grundzustand einen Regimewechsel zwischen einer ABCE-dominierten und einer AABB-dominierten neutralen Ordnung. Der Umschlag liegt numerisch bei
		\[
		\alpha_c\approx -0.72.
		\]
	\end{proposition}
	
	\begin{proof}
		Bei \(\alpha=-0.80\) ist
		\[
		\wt_{ABCE}(E_0)-\wt_{AABB}(E_0)>0,
		\]
		bei \(\alpha=-0.70\) dagegen
		\[
		\wt_{ABCE}(E_0)-\wt_{AABB}(E_0)<0.
		\]
		Der Vorzeichenwechsel lokalisiert den Umschlag zwischen diesen beiden Werten. Lineare Interpolation ergibt näherungsweise \(\alpha_c\approx -0.724\).
	\end{proof}
	
	\section{Schluss}
	
	Das diskrete Vierlingsmodell besitzt bereits in seiner einfachsten Form mehrere strukturelle Merkmale, die für das Bamberger Modell relevant sind:
	
	\begin{enumerate}[label=(\roman*)]
		\item einen globalen neutralen Sektor mit nichttrivialer interner Organisation,
		\item eine natürliche Viererorbit- und Shift-Struktur,
		\item diskrete Holonomie in den zyklischen Sektoren,
		\item konkurrierende Formen globaler Neutralität,
		\item einen numerisch lokalisierbaren Ordnungswechsel im Grundzustand.
	\end{enumerate}
	
	Der wichtigste Befund lautet: Der Grundzustand ist nicht durch triviale Leere, sondern durch strukturierte Neutralität bestimmt. Diese strukturierte Neutralität tritt in mindestens zwei Formen auf, nämlich als paarig kompensierte AABB-Ordnung und als voll gemischte ABCE-Ordnung.
	
	\section*{Ausblick}
	
	Naheliegende Fortsetzungen sind:
	\begin{enumerate}[label=(\arabic*)]
		\item ein zweidimensionaler Scan in \((\alpha,\lambda)\),
		\item die Analyse zusätzlicher Symmetrien und Entartungen,
		\item die Einführung lokaler Observablen und Korrelationsmaße,
		\item verallgemeinerte Holonomie-Terme,
		\item eine präzisere Zuordnung arithmetischer Vierlingsobjekte zu Signaturvierlingen.
	\end{enumerate}
	
	\begin{thebibliography}{9}
		
		\bibitem{AB}
		Y. Aharonov, D. Bohm,
		\textit{Significance of Electromagnetic Potentials in the Quantum Theory},
		Phys. Rev. \textbf{115} (1959), 485--491.
		
		\bibitem{vK}
		K. v. Klitzing, G. Dorda, M. Pepper,
		\textit{New Method for High-Accuracy Determination of the Fine-Structure Constant Based on Quantized Hall Resistance},
		Phys. Rev. Lett. \textbf{45} (1980), 494--497.
		
		\bibitem{ConwaySmith}
		J. H. Conway, D. A. Smith,
		\textit{On Quaternions and Octonions},
		A K Peters, 2003.
		
	\end{thebibliography}
	
\end{document}